%% notes-tex/026-kinetics-thermodynamics/sections/6-limits.tex
%% [[experiments.026-metabolism-flux.kinetic-and-thermodynamic-parameter-datasets]]
%% https://github.com/Mjvolk3/torchcell/tree/main/notes-tex/026-kinetics-thermodynamics/sections/6-limits.tex

\section{What these tables cannot support\secstatus{todo}}
\label{sec:limits}

\subsection{A predicted turnover number is not a measured one\secstatus{todo}}

Every row in the $k_{cat}$ and $K_M$ builds is a model output. The tables record that in a
provenance column beside the value, because a number with no provenance becomes
indistinguishable from a measurement within two steps, and the coverage gate has to be able
to separate an experimental value, a predicted value and a filled default.

The size of the distinction is measurable. Two predictors trained on overlapping data,
applied to the same 7{,}351 pairs, disagree by a median factor of about three. Neither
model's own reported error is that large. The disagreement is the honest uncertainty on any
single value, and it is wider than any published accuracy figure would suggest, because
those figures are computed on held-out data drawn from the same corpus rather than on a
genome-scale application.

\subsection{Accuracy against the Open Enzyme Database is memorization, not generalization\secstatus{todo}}

Comparing predictions to the Open Enzyme Database looks like validation and is not. The
database aggregates BRENDA and Sabio-RK, which is what these models were trained on, so a
matched enzyme-substrate pair is almost certainly inside their training sets. Numbers from
that comparison are reported here explicitly as memorization checks.

The join also has to be on the \emph{pair}. Joining on enzyme alone compares a prediction
for one substrate against a measurement for another, and doing so moves one of these
correlations from 0.98 to 0.82 and another from 0.62 to 0.12. The looser join is the one
that is easy to write by accident.

\subsection{$K_M$ is built and deliberately not consumed\secstatus{todo}}

The flux layer does not use $K_M$. The parameter enters only through a saturation factor
$\eta_{\mathrm{sat}} = \prod_i c_i/(K_{M,i} + c_i)$, which needs intracellular metabolite
concentrations, and those are measured for 191 of 2{,}806 metabolites, 6.8 percent. Adding
$K_M$ today would add a parameter and a free variable per reaction with no anchor for
either.

The measured argument for adding it anyway is promiscuity: underground reactions show
roughly twice the $K_M$ at indistinguishable $k_{cat}$, so a $k_{cat}$-only model cannot
tell a promiscuous side activity from a native one. The measured argument against is that
the $K_M$ predictions here are the least reliable of the set, and treating them as constants
rather than as order-of-magnitude priors would put that unreliability directly into a
constraint.

\subsection{The thermodynamic gaps\secstatus{todo}}

Three, in order of how much they bind.

\textbf{Novel compounds need a licensed tool.} Adding a compound eQuilibrator has never
seen requires pKa estimates from ChemAxon's \texttt{cxcalc}. Without the license the
compound cache is read-only, so the betaxanthin intermediates that are not already in it
cannot be priced at all. This is the binding constraint on applying the method to
heterologous pathways, which is exactly where it is most wanted.

\textbf{Coverage fell.} Reaction coverage went from 77.7 percent to 42.5 percent, because a
reaction needs every participant resolved. The unresolved metabolites are dominated by
acyl-CoA isomers and combinatorial lipid species, which are chemically ordinary and merely
un-indexed under the names the model uses.

\textbf{Compartment pH is not validated.} It is implemented, it produces shifts too large
to believe, and the pH table is assumed rather than sourced. Until that is diagnosed the
uniform build is the one to use, which means transport thermodynamics is set to zero by
construction and the proton motive force is absent from the model.

\subsection{What the mirror records when a predictor cannot run\secstatus{todo}}

All six now run. Where a predictor could not be run at all, the mirror records
\texttt{runnable: false} with the specific blocker, so that a gap in coverage is never
mistaken for a modeling choice.

Boost\_KM carried that flag until its weights were read rather than its filenames, and the
correction is recorded in Section~\ref{sec:mirrors} rather than quietly dropped: the
repository does ship fitted regressors, and the obstacle was never a missing model. It was
that UniRep is a TensorFlow~1 mLSTM stepping one residue at a time on CPU, which on a first
attempt occupied roughly 85 cores. The ESM-1b variant Wu et al. actually used replaces that
stage with a transformer that embeds all 887 enzymes on one GPU in four minutes.
